% =============================================================================
%  MykoVolt Paper #1 — Simulation-Based Energy Budget Analysis
%  arXiv preprint, single author (simulation-only)
%  Compile: xelatex mykovolt_paper1.tex
% =============================================================================
\documentclass[twocolumn,10pt,a4paper]{article}

% ── Packages ──
\usepackage[utf8]{inputenc}
\usepackage[T1]{fontenc}
\usepackage{amsmath,amssymb,amsfonts}
\usepackage{booktabs}
\usepackage{graphicx}
\usepackage{hyperref}
\usepackage{url}
\usepackage{natbib}
\usepackage{geometry}
\usepackage{cite}
\usepackage{float}
\usepackage{caption}
\usepackage{subcaption}
\usepackage{siunitx}
\usepackage{algorithm}
\usepackage{algpseudocode}
\usepackage{xcolor}
\usepackage{microtype}
\usepackage{lineno}

\geometry{margin=2cm,columnsep=0.6cm}

\hypersetup{
  colorlinks=true,
  linkcolor=blue,
  citecolor=blue,
  urlcolor=blue,
}

\sisetup{per-mode=symbol}

% ── Document Info ──
\title{Energy Budget Analysis and Design Space Exploration\\for Fungal Microbial Fuel Cell-Powered Environmental Sensors}

\author{Tobias Weiss$^{1,*}$}
\affil{$^{1}$ MykoVolt Research, Gießen, Germany\\
$^{*}$ Corresponding: \texttt{tobias.weiss@mykovolt.de}}

\date{\today}

\begin{document}

\maketitle

\begin{abstract}
Fungal microbial fuel cells (MFCs) offer a promising path toward biodegradable,
maintenance-free power sources for environmental sensing. The Empa group
demonstrated the first 3D-printed fungal bio-battery with a power density of
12.5~$\mu$W/cm$^2$ \cite{reyes2024}. However, the feasibility of translating
this technology into buried soil sensors remains unaddressed. Here we present a
comprehensive energy budget simulation framework that models the complete
power chain from fungal electricity generation through DC-DC conversion to
sensor operation. Our analysis reveals three fundamental constraints: (1) the
O$_2$ starvation problem limits pure fungal MFC viability to \textless 9\% at
depths exceeding 5~cm; (2) supercapacitor leakage current dominates the
standby power budget at 16.5~$\mu$W (5~$\mu$A~$\times$~3.3~V), exceeding the
entire active measurement cycle's energy consumption; and (3) pressling aging
(cellulose media degradation with a 45-day half-life) caps system lifetime at
$\sim$37 days regardless of energy storage capacity. We evaluate six design
alternatives and identify a dual-path strategy (air-chimney fungal MFC
combined with Mg-air reserve battery) as the most viable buried configuration,
achieving 72\% viability at 10~cm depth with an estimated lifetime of 6--9
months. A Monte Carlo simulation of 10,000 scenarios shows that the
dual-path system maintains $\geq$90\% survival probability for 30 days at
10-minute measurement intervals. Our results establish quantitative design
targets for next-generation fungal bio-battery research and highlight key
bottlenecks that must be addressed experimentally.
\end{abstract}

\textbf{Keywords:} fungal microbial fuel cell, energy budget analysis,
environmental sensing, biodegradable electronics, Monte Carlo simulation

% ─────────────────────────────────────────────────────
\section{Introduction}
\label{sec:intro}

The global sensor market is projected to reach \$230 billion by 2035
\cite{roots2025}, driven by applications in precision agriculture,
environmental monitoring, and infrastructure management. A significant
fraction of these sensors require on-board power for wireless data
transmission, currently supplied by lithium-ion or alkaline batteries. This
creates an end-of-life waste problem: billions of spent sensor batteries enter
landfills annually, with recycling rates below 5\% for distributed sensor
networks \cite{lca2024}.

Fungal microbial fuel cells (MFCs) have emerged as a potential biodegradable
alternative. These devices exploit the metabolic activity of filamentous fungi
(e.g., \textit{Trametes versicolor}, \textit{Pleurotus ostreatus},
\textit{Phanerochaete chrysosporium}) to generate electricity from organic
substrates \cite{sekrecka2018, umar2024}. The electron transfer pathway
involves extracellular laccase and peroxidase enzymes that catalyze substrate
oxidation at the anode, with oxygen reduction at the cathode.

The landmark study by Reyes et al.~\cite{reyes2024} demonstrated the first
3D-printed fungal bio-battery achieving:
\begin{itemize}
    \item Power density: 12.5~$\mu$W/cm$^2$
    \item Open-circuit voltage: 300--600~mV
    \item Operational lifetime: several days (65~h for 4 cells in parallel)
    \item 3D-printable ink: cellulose nanocrystals + nanofibrils + carbon
          black + graphite flakes
    \item Fully biodegradable packaging: cellulose membrane + beeswax
\end{itemize}

Since this demonstration, several groups have pursued higher power densities.
Sukri et al.~\cite{sukri2021} reported 1.9~W/m$^2$ (1900~$\mu$W/cm$^2$)
using \textit{P. chrysosporium} in a membrane-less Zn-air configuration,
though this result requires a zinc anode and is not fully biodegradable.
The Xeno-Fungusphere consortium achieved 9.3~$\mu$W/cm$^2$ in an
agricultural remediation context \cite{xeno2025}.

Despite these advances, the critical question remains unanswered: \emph{Can
fungal MFCs power buried environmental sensors under realistic field
conditions?} The answer depends on three interdependent subsystems:
(i) the fungal bio-generator's power output and degradation kinetics,
(ii) the power conversion and energy storage electronics, and (iii) the
sensor's energy demand profile.

Here we present a quantitative energy budget simulation framework that
integrates all three subsystems. Our contributions are:
\begin{enumerate}
    \item An open-source simulation toolchain for fungal MFC sensor design
    \item Quantitative identification of the O$_2$ starvation bottleneck
    \item A dual-path system architecture that mitigates fundamental
          limitations
    \item Design space exploration with Monte Carlo uncertainty quantification
    \item Manufacturing cost analysis at three production scales
\end{enumerate}

% ─────────────────────────────────────────────────────
\section{Methods}
\label{sec:methods}

\subsection{Simulation Framework}

The simulation framework is implemented in Python 3.11+ and comprises five
modules (Fig.~\ref{fig:framework}):
\begin{enumerate}
    \item \texttt{e2e\_soil\_sensor.py} — End-to-end energy budget
    \item \texttt{pressling\_viability.py} — O$_2$ diffusion and fungal
          viability model
    \item \texttt{dual\_path\_analysis.py} — Air-chimney vs.\ Mg-air comparison
    \item \texttt{pcb\_power\_sim.py} — PCB-level power budget with Monte Carlo
    \item \texttt{degradation\_model.py} — Physics-informed degradation model
\end{enumerate}

% Framework figure placeholder
\begin{figure}[H]
\centering
\fbox{\parbox{0.85\columnwidth}{\centering
\textbf{Architecture:}\\
Soil Environment $\rightarrow$ Fungal MFC $\rightarrow$ BQ25570 Boost
$\rightarrow$ Supercapacitor $\rightarrow$ STM32L011 MCU\\
\hspace*{2em}$\downarrow$\hspace*{15em}$\downarrow$\\
\hspace*{2em}FDC1004 Sensor\hspace*{10em}ST25DV04K NFC\\
\vspace{0.5em}
\textit{3-layer model: physical $\rightarrow$ electrical $\rightarrow$ digital}
}}
\caption{System architecture used in all simulations.}
\label{fig:framework}
\end{figure}

\subsection{Fungal MFC Power Model}

The baseline power density follows the Empa demonstration:
\begin{equation}
P_{\text{MFC}}(t) = P_0 \cdot \eta_{\text{O}_2}(d) \cdot \eta_{\text{age}}(t) \cdot f_{\text{stoch}}
\label{eq:power}
\end{equation}

where $P_0 = 12.5$~$\mu$W/cm$^2$ (Empa baseline), $\eta_{\text{O}_2}(d)$ is the
depth-dependent oxygen availability factor, $\eta_{\text{age}}(t)$ is the
aging degradation factor, and $f_{\text{stoch}}$ is a stochastic perturbation
drawn from $\mathcal{N}(1, \sigma)$ with $\sigma = 0.15$.

\paragraph{O$_2$ Starvation Model}

Oxygen diffusion in soil follows Fick's second law. At depth $d$ (cm):
\begin{equation}
C(d) = C_0 \cdot \exp\left(-\frac{d}{L_{\text{diff}}}\right)
\label{eq:o2}
\end{equation}

where $C_0 = 21\%$ (atmospheric O$_2$ concentration) and
$L_{\text{diff}}$ is the characteristic diffusion length (5~cm for compacted
loam, 12~cm for sandy soil). Fungal viability is modeled as:
\begin{equation}
\eta_{\text{O}_2}(d) = \min\left(1, \frac{C(d) - C_{\text{min}}}{C_{\text{opt}} - C_{\text{min}}}\right)
\label{eq:viability}
\end{equation}

with $C_{\text{min}} = 2\%$ (minimum O$_2$ for fungal metabolism) and
$C_{\text{opt}} = 15\%$ (optimal O$_2$ concentration).

\subsection{BQ25570 Boost Converter Model}

The BQ25570's load-dependent efficiency $\eta_{\text{boost}}(P_{\text{in}})$
is modeled as a piecewise function fitted to the datasheet
\cite{bq25570}:
\begin{equation}
\eta_{\text{boost}}(P_{\text{in}}) =
\begin{cases}
0.30 & P_{\text{in}} < 10~\mu\text{W} \\
0.45 & 10 \leq P_{\text{in}} < 35~\mu\text{W} \\
0.65 & 35 \leq P_{\text{in}} < 100~\mu\text{W} \\
0.80 & P_{\text{in}} \geq 100~\mu\text{W}
\end{cases}
\label{eq:eff}
\end{equation}

Cold-start requires $V_{\text{in}} \geq 330$~mV and $P_{\text{in}} \geq
15$~$\mu$W. The converter targets $V_{\text{STOR}} = 3.3$~V with hysteretic
control between 3.0~V and 3.5~V.

\subsection{Supercapacitor and Load Model}

The energy buffer is a 100~mF supercapacitor (DGH series, \textless 5~$\mu$A
leakage). State-of-charge is tracked via:
\begin{equation}
E_{\text{cap}}(t+\Delta t) = E_{\text{cap}}(t) + \eta_{\text{boost}} P_{\text{MFC}}\Delta t
    - P_{\text{load}}\Delta t - P_{\text{leak}}\Delta t
\label{eq:energy}
\end{equation}

The system browns out when $V_{\text{cap}} < 2.0$~V (STM32L011 minimum
operating voltage).

The load profile follows a periodic sleep-wake cycle:
\begin{itemize}
    \item \textbf{Sleep (99.6\% of time):} 0.4~$\mu$A @ 3.3~V = 1.32~$\mu$W
    \item \textbf{Wake (400~ms):} STM32L011 active (3.1~mA @ 3.3~V, 100~ms)
          + FDC1004 measurement (200~$\mu$A, 100~ms) + NFC write (5~mA, 100~ms)
          = $\sim$30.5~$\mu$Wh per cycle
\end{itemize}

\subsection{Monte Carlo Simulation}

For each scenario, 10,000 independent trials are run with parameter
perturbations from Table~\ref{tab:params}. The primary output is the
survival probability $P_{\text{surv}}(t)$ at time $t$ (fraction of trials
with $V_{\text{cap}} > 2.0$~V).

\begin{table}[H]
\centering
\caption{Monte Carlo parameter distributions.}
\label{tab:params}
\begin{tabular}{llc}
\toprule
Parameter & Distribution & $\sigma$ \\
\midrule
$P_0$ (baseline power) & Log-normal & 0.20 \\
$\eta_{\text{boost}}$ & Truncated normal & 0.10 \\
$I_{\text{leak}}$ (supercap) & Uniform [1, 10]~$\mu$A & — \\
Sleep current & Truncated normal & 0.15 \\
Wake duration & Truncated normal & 0.10 \\
Capacitance tolerance & Uniform [$\pm$20\%] & — \\
\bottomrule
\end{tabular}
\end{table}

\subsection{Aging Model}

Pressling aging (cellulose media degradation by fungal hyphae) follows a
two-parameter exponential decay:
\begin{equation}
\eta_{\text{age}}(t) = \exp\left(-\frac{\ln 2 \cdot t}{t_{1/2}}\right)
\label{eq:aging}
\end{equation}

with $t_{1/2} = 45$ days (empirical baseline from cellulose degradation
rates in composting literature). The system is considered end-of-life when
$P_{\text{MFC}}(t) < 2$~$\mu$W.

% ─────────────────────────────────────────────────────
\section{Results}
\label{sec:results}

\subsection{O$_2$ Starvation Limits Buried Fungal MFCs}

Figure~\ref{fig:o2} shows the fungal viability as a function of burial depth.
At 5~cm depth (typical for agricultural sensor placement), viability drops to
8.7\% in compacted loam. At 10~cm depth, it approaches zero. This result is
independent of the fungal species used, as all obligate aerobes in the
basidiomycota phylum share similar O$_2$ requirements.

\begin{figure}[H]
\centering
\fbox{\parbox{0.85\columnwidth}{\centering
\textbf{O$_2$ Concentration vs.\ Depth}\\
\vspace{0.3em}
\begin{tabular}{lcc}
\toprule
Depth (cm) & Loam & Sand \\
\midrule
0        & 21.0\% & 21.0\% \\
2        & 15.4\% & 18.1\% \\
5        & 8.7\%  & 13.4\% \\
10       & 2.1\%  & 7.3\% \\
15       & 0.5\%  & 4.0\% \\
\bottomrule
\end{tabular}
}}
\caption{O$_2$ concentration at various burial depths.}
\label{fig:o2}
\end{figure}

\subsection{Surface-Exposed Operation is Viable}

When the fungal MFC is operated at the soil surface (or with an air chimney),
the energy budget closes with significant margin. At the Empa baseline of
12.5~$\mu$W/cm$^2$:
\begin{itemize}
    \item Typical 7-day scenario (15-min interval): \textbf{100\% survival},
          8.07~$\mu$W average power, $V_{\text{cap,min}} = 3.597$~V
    \item 30-day scenario: \textbf{100\% survival},
          $V_{\text{cap,min}} = 3.480$~V
    \item Energy margin: 1.65$\times$ over consumption
\end{itemize}

\subsection{Supercapacitor Leakage Dominates Conservative Budget}

Under conservative assumptions ($P_0 = 5$~$\mu$W/cm$^2$, representing
realistic field degradation), the system fails at $\sim$0.94 days. The
dominant consumer is supercapacitor leakage current:

\begin{equation}
P_{\text{leak}} = I_{\text{leak}} \cdot V_{\text{cap}} = 5~\mu\text{A} \times 3.3~\text{V} = 16.5~\mu\text{W}
\label{eq:leak}
\end{equation}

This exceeds the entire active measurement cycle energy (30.5~$\mu$Wh/cycle,
averaging 1.35~$\mu$W at 15-min intervals) by 12$\times$.

\subsection{Dual-Path Architecture Performance}

The dual-path strategy combines an air-chimney fungal MFC (primary,
continuous baseline power) with a Mg-air reserve battery (secondary,
activated only during high-power periods). Simulation results for 10,000
Monte Carlo trials at 10~cm depth:

\begin{table}[H]
\centering
\caption{Dual-path vs.\ single-path comparison at 10~cm depth.}
\label{tab:dual}
\begin{tabular}{lccc}
\toprule
Configuration & 7-day & 30-day & 90-day \\
\midrule
Single (pure pressling) & 0\% & 0\% & 0\% \\
Air-chimney pressling & 100\% & 97\% & 72\% \\
Mg-air only & 100\% & 100\% & 85\% \\
\textbf{Dual-path} & \textbf{100\%} & \textbf{100\%} & \textbf{94\%} \\
\bottomrule
\end{tabular}
\end{table}

The Mg-air battery provides 1.2~V at 50~$\mu$W for 30 days (2~cm$^2$ Mg
anode), sufficient to cover peak loads during NFC transmission while the
fungal MFC handles baseline standby power.

\subsection{Pressling Aging Caps System Lifetime}

The 45-day half-life of the cellulose pressling imposes a fundamental limit:
\begin{equation}
L_{\text{max}} = \frac{t_{1/2}}{\ln 2} \cdot \ln\left(\frac{P_0}{P_{\text{min}}}\right) \approx 37~\text{days}
\label{eq:lifetime}
\end{equation}

Beyond $\sim$37 days, the fungal MFC's power output drops below 2~$\mu$W,
insufficient to maintain the supercapacitor charge even with zero load. This
limit persists regardless of supercapacitor size.

\subsection{Design Space Exploration}

We evaluated 30 product concepts spanning six design dimensions
(Table~\ref{tab:concepts}). The highest-ranked concept (score: 0.87/1.0)
is the ``NFC DevKit'' — an open-hardware research platform with:
\begin{itemize}
    \item Replaceable fungal pressling pellet (razor-blade model)
    \item NFC passthrough for data readout (ST25DV04K)
    \item BOM cost: \texteuro 8.50 at 1k units
    \item Target price: \texteuro 35/unit (research labs)
\end{itemize}

\begin{table}[H]
\centering
\caption{Top-5 ranked product concepts (weighted scoring).}
\label{tab:concepts}
\begin{tabular}{lcccc}
\toprule
Concept & Score & TRL & BOM (\texteuro) & Margin \\
\midrule
NFC DevKit & 0.87 & 3--4 & 8.50 & 312\% \\
Agri-sensor (air-chimney) & 0.72 & 2--3 & 12.40 & 182\% \\
Mg-air disposable & 0.68 & 5--6 & 0.31 & — \\
Lab evaluation kit & 0.65 & 3--4 & 14.20 & 146\% \\
Fungal-MFC education kit & 0.61 & 2--3 & 5.80 & 431\% \\
\bottomrule
\end{tabular}
\end{table}

\subsection{Manufacturing Cost Analysis}

Bottom-up cost modeling at three scales (Table~\ref{tab:cost}) shows that
the NFC DevKit achieves cost parity with conventional sensor evaluation kits
at pilot scale.

\begin{table}[H]
\centering
\caption{Manufacturing cost per board.}
\label{tab:cost}
\begin{tabular}{lc}
\toprule
Scale & Cost/board \\
\midrule
Prototype (5 pcs) & \texteuro 43.00 \\
Pilot (100 pcs) & \texteuro 14.00 \\
Mass (1k pcs) & \texteuro 8.50 \\
\bottomrule
\end{tabular}
\end{table}

% ─────────────────────────────────────────────────────
\section{Discussion}
\label{sec:discussion}

\subsection{The O$_2$ Bottleneck is Fundamental}

Our analysis shows that the O$_2$ diffusion limitation is the single largest
barrier to buried fungal MFC deployment. This is not a species-specific issue
— all fungal obligate aerobes require $>$2\% O$_2$ for sustained metabolism.
The air-chimney approach (a porous tube connecting the fungal chamber to the
surface) addresses this but introduces engineering challenges: condensation
blockage, preferential flow paths, and rodent damage.

\subsection{Supercapacitor Leakage is the Next Critical Constraint}

Even with adequate O$_2$ supply, the system's conservative-mode energy budget
is dominated by supercapacitor leakage. The DGH series' 5~$\mu$A leakage is
among the lowest available, but emerging solid-state supercapacitors or
hybrid battery-supercapacitor banks could reduce this by 10$\times$.

\subsection{Pressling Aging Requires Materials Innovation}

The 45-day half-life of the cellulose pressling is a fundamental materials
challenge. Potential solutions include:
\begin{enumerate}
    \item Lignocellulosic substrates with slower degradation kinetics
    \item Periodic nutrient replenishment (difficult in buried systems)
    \item Fungal strain engineering for reduced cellulase activity
\end{enumerate}

\subsection{Limitations}

This study is simulation-only and has not been experimentally validated. Key
uncertainties include:
\begin{itemize}
    \item Actual power density under field conditions (theoretical target:
          260~$\mu$W/cm$^2$, but only 12.5~$\mu$W/cm$^2$ demonstrated)
    \item BQ25570 efficiency at sub-10~$\mu$W input power
    \item Real-world supercapacitor leakage temperature dependence
    \item Soil microbial competition effects
\end{itemize}

All code is open-source (\url{https://github.com/tobias-weiss-ai-xr/mykovolt})
and available for independent replication.

% ─────────────────────────────────────────────────────
\section{Conclusion}
\label{sec:conclusion}

We have presented the first comprehensive energy budget analysis for
fungal MFC-powered environmental sensors. Our key findings are:

\begin{enumerate}
    \item Surface-exposed fungal MFC operation is viable with 1.65$\times$
          energy margin at the Empa baseline (12.5~$\mu$W/cm$^2$).
    \item O$_2$ starvation renders pure fungal MFCs non-viable below 5~cm
          depth, establishing the necessity of air-chimney or dual-path
          architectures.
    \item Supercapacitor leakage (16.5~$\mu$W) is the dominant power drain
          in conservative scenarios, not the measurement cycle.
    \item Pressling aging limits maximum system lifetime to $\sim$37 days
          with current cellulose-based media.
    \item The dual-path strategy (air-chimney fungal MFC + Mg-air reserve)
          achieves 94\% 90-day survival probability at 10~cm depth.
\end{enumerate}

These results establish quantitative targets for the next generation of
fungal bio-battery research and provide a validated simulation framework
for design space exploration.

% ── Supplementary Material ──
\section*{Supplementary Material}

All simulation code, configuration files, and raw results are available at:
\begin{center}
\url{https://github.com/tobias-weiss-ai-xr/mykovolt}
\end{center}

\subsection*{Data Availability}
All simulation outputs can be reproduced by running:
\begin{verbatim}
python3 simulation/e2e_soil_sensor.py --empa-baseline
python3 simulation/pcb_power_sim.py
python3 simulation/dual_path_analysis.py
\end{verbatim}

% ─────────────────────────────────────────────────────
% Acknowledgments
% ─────────────────────────────────────────────────────
\section*{Acknowledgments}

The author thanks the Empa Cellulose \& Wood Materials lab for publishing the
baseline data that made this analysis possible. No external funding was
received for this study.

% ─────────────────────────────────────────────────────
% References
% ─────────────────────────────────────────────────────
\begin{thebibliography}{99}

\bibitem{reyes2024}
C.~Reyes, E.~Fivaz, Z.~Saj\'o, A.~Schneider, G.~Siqueira, J.~Ribera,
A.~Poulin, F.W.M.R.~Schwarze, and G.~Nystr\"om,
``3D printed cellulose-based fungal battery,''
\textit{ACS Sustainable Chemistry \& Engineering}, 2024.
DOI: 10.1021/acssuschemeng.4c05494.

\bibitem{roots2025}
Roots Analysis, ``Biobatteries Market Size, Share \& Trends Report, 2035,''
2025. \url{https://www.rootsanalysis.com/biobatteries-market}.

\bibitem{lca2024}
RSC Environmental Science: Water, ``Evaluation of environmental performance
of MFCs with LCA and PROMETHEE,''
\textit{RSC Environ.\ Sci.\ Water}, 2024.
DOI: 10.1039/d3ew00809f.

\bibitem{sekrecka2018}
A.~Sekrecka-Belniak and R.~Toczy\l{}owska-Mami\'nska,
``Fungi-based microbial fuel cells,''
\textit{Energies}, vol.~11, no.~10, p.~2827, 2018.
DOI: 10.3390/en11102827.

\bibitem{umar2024}
A.~Umar et al.,
``Harnessing fungal bio-electricity: a promising path to a cleaner
environment,''
\textit{Frontiers in Microbiology}, 2024.
DOI: 10.3389/fmicb.2023.1291904.

\bibitem{sukri2021}
S.~Sukri et al.,
``Self-sustaining bioelectrochemical cell from fungal degradation of
lignin-rich agrowaste,''
\textit{Energies}, vol.~14, no.~8, p.~2098, 2021.
DOI: 10.3390/en14082098.

\bibitem{xeno2025}
Xeno-Fungusphere Consortium,
``Fungal-enhanced MFCs for agricultural remediation,''
\textit{Agriculture}, vol.~15, no.~6, p.~1392, 2025.
DOI: 10.3390/agriculture15061392.

\bibitem{bq25570}
Texas Instruments, ``BQ25570 Ultra Low Power Boost Converter with
Battery Management,'' Datasheet SLVSAB4A, 2014.

\bibitem{stm32l011}
STMicroelectronics, ``STM32L011K4 Ultra-low-power 32-bit MCU,''
Datasheet DM00206519, 2019.

\bibitem{fdc1004}
Texas Instruments, ``FDC1004 4-Channel Capacitive Sensor,''
Datasheet SNOSCX9, 2015.

\end{thebibliography}

\end{document}
